\documentclass[12pt,a4paper]{article}

\usepackage[utf8]{inputenc}
\usepackage[T1]{fontenc}
\usepackage[spanish]{babel}
\usepackage{geometry}
\usepackage{xcolor}
\usepackage{listings}
\usepackage{hyperref}
\usepackage{array}
\usepackage{longtable}

\geometry{margin=2.5cm}

\hypersetup{
    colorlinks=true,
    linkcolor=black,
    urlcolor=blue
}

\lstdefinestyle{codigo}{
    basicstyle=\ttfamily\small,
    breaklines=true,
    frame=single,
    numbers=left,
    numberstyle=\tiny,
    keywordstyle=\color{blue},
    commentstyle=\color{gray},
    stringstyle=\color{red},
    tabsize=4
}

\lstset{style=codigo}

\begin{document}

\begin{titlepage}
    \centering
    {\Large Universidad Nacional San Antonio Abad del Cusco\par}
    \vspace{0.3cm}
    {\large Escuela Profesional de Ingenieria Informatica y de Sistemas\par}
    \vspace{1.5cm}
    {\Huge\bfseries Laboratorio Nro. 06\par}
    \vspace{0.4cm}
    {\Large Analizador Sintactico LR con Flex y Bison\par}
    \vspace{1.5cm}
    \begin{tabular}{>{\bfseries}ll}
        Curso: & IF454AIN - Compiladores \\
        Profesor: & Victor Dario Sosa Jauregui \\
        Estudiante: & Efrain Vitorino Maric \\
        Codigo: & 160337 \\
        Repositorio: & \url{https://github.com/lovito99/LaboCompiladores.git} \\
    \end{tabular}
    \vfill
    {\large Cusco, 2026\par}
\end{titlepage}

\tableofcontents
\newpage

\section{Introduccion}

El presente informe describe la implementacion del Laboratorio 6, que consiste en un analizador sintactico LR construido con Flex y Bison. El programa reconoce expresiones aritmeticas, aplica jerarquia de operadores, evalua cada expresion y reporta errores lexicos, sintacticos y matematicos.

Las mejoras incluidas son el operador modulo o resto, el operador potencia, la aceptacion de numeros decimales positivos usando \texttt{atof}, y el manejo de errores para division por cero y cero elevado a cero.

\section{Objetivos}

\begin{itemize}
    \item Reconocer expresiones aritmeticas mediante un analizador lexico y sintactico.
    \item Implementar los operadores suma, resta, multiplicacion, division, modulo y potencia.
    \item Verificar la jerarquia y asociatividad de operadores.
    \item Aceptar numeros enteros y decimales positivos.
    \item Reportar errores matematicos para division por cero, modulo por cero y cero potencia cero.
\end{itemize}

\section{Archivos del proyecto}

\begin{longtable}{|p{4cm}|p{10cm}|}
\hline
\textbf{Archivo} & \textbf{Descripcion} \\
\hline
\texttt{lexer.l} & Analizador lexico. Reconoce numeros, operadores, parentesis y saltos de linea. \\
\hline
\texttt{analizador.y} & Analizador sintactico. Define gramatica, precedencia, acciones semanticas y errores. \\
\hline
\texttt{Makefile} & Automatiza la compilacion, ejecucion y limpieza del proyecto. \\
\hline
\texttt{informe/} & Carpeta que contiene el documento LaTeX y el PDF del informe. \\
\hline
\end{longtable}

\section{Fragmentos de codigo}

\subsection{Reconocimiento de numeros decimales positivos}

El analizador lexico reconoce enteros positivos, decimales con parte entera y decimales que inician con punto. El texto reconocido se convierte a \texttt{double} usando \texttt{atof}.

\begin{lstlisting}[language=C]
[0-9]+(\.[0-9]+)?|\.[0-9]+ {
    yylval.valor = atof(yytext);
    return NUM;
}
\end{lstlisting}

\subsection{Precedencia y asociatividad}

La jerarquia se define directamente en Bison. Las declaraciones escritas despues tienen mayor prioridad.

\begin{lstlisting}[language=C]
%left '+' '-'          /* menor prioridad */
%left '*' '/' '%'      /* prioridad media */
%right UMINUS          /* signo menos unario */
%right POTENCIA        /* mayor prioridad */
\end{lstlisting}

Con esta configuracion, la potencia tiene mayor prioridad que el producto, division, modulo, suma y resta. Ademas, \texttt{POTENCIA} es asociativa por la derecha, por lo que \texttt{2**3**2} se interpreta como \texttt{2**(3**2)}.

\subsection{Modulo o resto}

El operador modulo se evalua usando \texttt{fmod}, funcion de la libreria matematica.

\begin{lstlisting}[language=C]
| expr '%' expr
  {
      if ($3 == 0) {
          errores_sintacticos++;
          error_matematico = 1;
          fprintf(stderr, "Linea %d: Error matematico. Modulo por cero.\n",
                  @3.first_line);
          $$ = 0;
      } else {
          $$ = fmod($1, $3);
      }
  }
\end{lstlisting}

\subsection{Potencia}

La potencia se evalua con \texttt{pow}. Se agrega una validacion especial para el caso \texttt{0**0}.

\begin{lstlisting}[language=C]
| expr POTENCIA expr
  {
      if ($1 == 0 && $3 == 0) {
          errores_sintacticos++;
          error_matematico = 1;
          fprintf(stderr,
                  "Linea %d: Error matematico. Cero elevado a cero no esta definido.\n",
                  @2.first_line);
          $$ = 0;
      } else {
          $$ = pow($1, $3);
      }
  }
\end{lstlisting}

\subsection{Division por cero}

La division valida el divisor antes de calcular el resultado.

\begin{lstlisting}[language=C]
| expr '/' expr
  {
      if ($3 == 0) {
          errores_sintacticos++;
          error_matematico = 1;
          fprintf(stderr, "Linea %d: Error matematico. Division por cero.\n",
                  @3.first_line);
          $$ = 0;
      } else {
          $$ = $1 / $3;
      }
  }
\end{lstlisting}

\section{Compilacion}

La compilacion se realiza desde la carpeta \texttt{laboratorio-6}.

\begin{lstlisting}[language=bash]
make clean
make all
\end{lstlisting}

Salida obtenida:

\begin{lstlisting}[language=bash]
rm -f analizador analizador.tab.c analizador.tab.h lexer.yy.c analizador.output
bison -d analizador.y -o analizador.tab.c
flex -o lexer.yy.c lexer.l
gcc -Wall -Wextra -g analizador.tab.c lexer.yy.c -o analizador -lm
rm -f analizador.tab.c analizador.tab.h lexer.yy.c
\end{lstlisting}

\section{Pruebas de ejecucion}

\subsection{Pruebas correctas}

Entrada:

\begin{lstlisting}
10%3
10.5%4
2**3**2
-2**2
(-2)**2
1.5+2.25
\end{lstlisting}

Salida:

\begin{lstlisting}
Resultado = 1
Resultado = 2.5
Resultado = 512
Resultado = -4
Resultado = 4
Resultado = 3.75
\end{lstlisting}

\subsection{Verificacion de jerarquia}

\begin{longtable}{|p{4cm}|p{6cm}|p{3cm}|}
\hline
\textbf{Expresion} & \textbf{Interpretacion} & \textbf{Resultado} \\
\hline
\texttt{2**3**2} & \texttt{2**(3**2)} & \texttt{512} \\
\hline
\texttt{-2**2} & \texttt{-(2**2)} & \texttt{-4} \\
\hline
\texttt{(-2)**2} & Parentesis cambian la prioridad & \texttt{4} \\
\hline
\texttt{1.5+2.25} & Suma de decimales & \texttt{3.75} \\
\hline
\end{longtable}

\section{Informe de errores}

Entrada:

\begin{lstlisting}
1/0
0**0
10%0
\end{lstlisting}

Salida:

\begin{lstlisting}
Linea 1: Error matematico. Division por cero.
Linea 2: Error matematico. Cero elevado a cero no esta definido.
Linea 3: Error matematico. Modulo por cero.
\end{lstlisting}

Estos errores no imprimen resultado numerico para la expresion incorrecta. El programa marca el error matematico y continua con la siguiente linea de entrada.

\section{Conclusiones}

Se implemento el analizador sintactico LR con soporte para numeros decimales positivos, modulo y potencia. La jerarquia de operadores se controla mediante declaraciones de precedencia de Bison. Tambien se agrego manejo de errores matematicos para evitar resultados invalidos en division por cero, modulo por cero y cero elevado a cero.

\section{Repositorio}

El codigo fuente del laboratorio se encuentra disponible en:

\begin{center}
\url{https://github.com/lovito99/LaboCompiladores.git}
\end{center}

\end{document}
